\section{Related Literature}
\label{sec:literature}

This paper makes four moves on existing screening and bid-rigging
literatures.

\paragraph{Bid-coordination tests require ex-post partitioning.}
\citet{bajari2003deciding} develop the canonical tests for
collusion---exchangeability of bid functions and conditional
independence of bid residuals---but the researcher must first
partition bidders into suspected colluders and a competitive
fringe, a partition the tests themselves cannot deliver. The
literature that followed inherits this constraint: bid rotation
\citep{porter1993detection, porter1999ohio}, collusion in timber
auctions \citep{baldwin1997bidder}, bidder groups from co-bidding
patterns \citep{conley2016detecting}, and collusion dynamics in
long-lived cartels \citep{clark2021collusion,
schurter2020identification} all classify firms \emph{after}
collusion is established or assumed. \citet{conley2016detecting}
come closest in data parsimony but detect group structure rather
than flag individual firms. Our contribution is the partitioning
step: a data-driven, \emph{ex ante} firm-level classifier that
produces the input these subsequent tests need.

\paragraph{Proactive screens require bid microdata that does not
always exist.}
\citet{imhof2019detecting}, \citet{imhof2018screening},
\citet{huber2019machine}, and \citet{wallimann2023machine} train
machine-learning classifiers on tender-level bid statistics
(within-tender CV, kurtosis, normalized spread) and achieve high
accuracy against known cartels; \citet{abrantesmetz2006a}
develop the variance screen; \citet{chassang2022robust} provide
a robust theoretical foundation; \citet{harrington2008detecting}
surveys. All four operate at the \emph{tender} level and require
granular bid microdata---in particular, the per-bidder bid
amounts. Many real procurement systems do not preserve this layer:
electronic platforms typically archive winning prices and
participant lists at the contract-record layer, not bid-by-bid.%
\footnote{Our Imhof comparison uses a coefficient-of-variation
median split, a coarser proxy for the full implementation; see
Section~\ref{sec:results_detection}.}
The gap between what cartel screens \emph{can} do (with full
microdata) and what oversight bodies \emph{do} have (winner +
participants) is the operational gap our screen addresses. The
FL classifier needs only contract-award records.

\paragraph{Cover-bidding identification has been retrospective.}
Cover bidders are documented in prosecuted cartels by
\citet{porter1999ohio}, \citet{pesendorfer2000study},
\citet{asker2010leniency}, and \citet{marshall2012economics}, and
in theoretical comparisons of auction formats by
\citet{athey2011comparing}. In every empirical case the cover
bidders are identified \emph{ex post}---through leniency
testimony, judicial documents, or post-conviction record review.
We invert this: by flagging firms whose participation patterns
are behaviorally coherent with cover bidding, the retrospective
observation becomes a prospective screening rule.

\paragraph{Procurement enforcement is hardest where data are
scarcest.}
The bid-rigging literature is dominated by developed-economy
settings (US, EU, OECD) where data infrastructure is mature.
Developing economies, where procurement spending shares of GDP
are larger and oversight resources thinner
\citep{sanchezgraells2019screening, oecd2019government}, are
both more vulnerable to bid rigging and less equipped to deploy
the existing detection tools. BEC's institutional setup---a
centralized electronic platform with 1{,}308 purchasing units,
4.5 million tender-items, 40 million bids, two auction formats,
a procurement-cartel adjudication portfolio at CADE---provides
the test bed for a screening rule designed to be deployable
under thin data infrastructure. The contribution generalizes
in principle wherever a procurement system records winners and
participants, regardless of bid-level archiving.
